%% Theory of Computing
%% File: v019a006.tex
%% Publication date: November 14, 2023
%% Authors: Linh Tran and Van Vu

%

%
%
%
%

%

%

%
%

%

%
%
%
%
%
%
%
%
	%
		%
		%
	%
	%
	%
%
%
%
%
%
%
%
%

%

\documentclass{toc}

%
%
%
\tocdetails{%
	title = {Reaching a Consensus on Random Networks: The Power of Few},
	%
	%
	%
	number_of_pages = {21}, 
	number_of_bibitems = {18},
	number_of_figures = {0},
	conference_version = {RANDOM20}, %
	%
	%
	%
	%
	%
	author = {Linh Tran and Van Vu},
	%
	%
	%
	%
	%
	authorlist = {Linh Tran, Van Vu},
	%
	%
	%
	%
	%
	%
	%
	%
	%
	%
	%
	%
	%
	%
	acmclassification = {G.3, G.2.2},
	%
	%
	amsclassification = {37E25, 05C57, 05C80, 82B43, 68R10, 60K35, 60K37},
	%
	%
	keywords = {random graphs, Erd\H{o}s-R\'enyi model, majority dynamics, consensus},
	%
	%
	%
	%
	%
}   %

%
%
%
%
%
%
%
%
%
%
%
%
%
%
%
%
%
%
%
%
%
%
%
%
%
%
%
%
%

%
%
\tocdetails{%
	volume = 19,      %
	number = 6,       %
	year = 2023,
	%
	%
	%
	%
	%
	%
	%
	%
	%
	specissue={\approxrandomCA},  %
	%
	received = {March 1, 2021},
	revised = {July 20, 2022},
	published = {November 14, 2023},
	%
	%
	%
	doi = {10.4086/toc.2023.v019a006},
	%
}   %

%
%
%
%
%
%
%

%
%
%
\usepackage{macros}

\newcommand{\err}{Erd\H{o}s--R{\'e}nyi}        

%
\usepackage{lineno}  %


\begin{document}
	
	\begin{frontmatter}%
		
		%
		%
		%
		
		%
		%
		
		%
		%
		
\title{Reaching a Consensus on Random Networks: The Power of Few%
%
\titlefootnote{A preliminary version of this paper appeared in  %
the \href{https://doi.org/10.4230/LIPIcs.APPROX/RANDOM.2020.20}%
%
{Proceedings of the 24th International Conference on Randomization and Computation   
(RANDOM'20)~\cite{conf-version}.}}}  %
%
			%
			%
				
				%
				%
				%
				%
				%
				%
				%
				%
				%
				%
				%
				%
				
				
				%
				%
				%
				%
				
				%
				
				\author[tran]{Linh Tran}
				\author[vu]{Van Vu}
				
				%
				%
					%
					%
				
				%
				%
				%
				%
				%
				%
				%
				%
				\begin{abstract}
					A community of $n$ individuals splits into two camps, Red and Blue. The individuals are connected by a social network, which influences their colors.
					%
					Every day   %
					each person changes %
					%
					%
					their  %
					color according to the majority of %
					their  %
					neighbors. Red (Blue) wins if everyone in the community becomes Red (Blue) 
					at some point.
					
					We study this process when the underlying network is the random 
					%
					\err\  %
					graph $G(n, p)$. With a balanced initial state ($n/2$ persons in each camp),
					it is clear that each color wins with the same probability.
					
					%
					Our study reveals that for any  constants $p$ and $\varepsilon$,
					there is a constant $c$ such that if one camp has at least $n/2 + c$ individuals at the initial state, then
					it wins with probability at least $1 - \varepsilon$. The surprising fact here is that $c$ \emph{does not} depend on $n$, the population of the community. When $p=1/2$ and $\varepsilon =.1$, one can set $c=5$, meaning one camp has $n/2 + 5$ members initially. In other words, it takes only $5$ extra people to win an election with overwhelming odds.
					We also generalize the result to $p = p_n = \text{o}(1)$ in a separate paper.
				\end{abstract}
				%
				
				%
				%
				%
				%
				%
				%
				%
				
				\iffalse %
				%
				%
				%
				%
				%
				%
				%
				%
				%
				%
				%
				%
				%
				%
				%
				%
				%
				%
				%
				%
				%
				
				\tocarxivcategory{!!!}
				
				\fi %
				
			\end{frontmatter}
			
			%
			%
			%
			%
			%
			%
			
			\section {Introduction} \label{sec:introduction}
			
			
			\subsection {The opinion exchange  dynamics}
			
			Building mathematical models to explain how collective opinions are formed is an important and interesting task (see \cite{mossel2017} for a survey on the topic, with examples from  various fields, economy, sociology, statistical physics, to mention a few).
			
			Obviously, our opinions are influenced by people around us, and this motivates the study of the following natural and simple model:
			A community of $n$ individuals splits into two camps,  Red and Blue, representing two competing opinions, which can be on any topic such as brand competition, politics, ethical issues, etc. The individuals are connected by a social network, which influences  their opinion on a daily basis (by some specific rule). We say that Red (respectively Blue) \emph{wins} if everyone in the community becomes Red (respectively Blue) at some point.
			
			
			We study  this process  when the underlying network is random. In this paper, we focus on the %
			\err\   %
			random graph $G(n, p)$, which is the most popular model of random graphs  \cite{Bollobasbook, Jansonbook}. We use the majority rule, which is a natural choice. When a new day comes, a vertex scans its neighbors' colors %
			from  %
			the previous day and adopts the dominant one. If there is a tie, it keeps its color.
			
			
			\begin{definition} The random graph $G(n,p)$ on $n\in \N$ vertices and density $p\in (0, 1)$ is obtained by putting an edge between any two vertices with probability $p$, independently.
			\end{definition}
			
			
			
			\subsection{Results}
			With a balanced initial state ($n/2$ persons in each camp), by symmetry, each color wins with the same  probability $q < 1/2$, regardless of $p$.   (Notice that there are graphs, such as the empty and complete graphs, on which no one wins.)
			Our study reveals  that for any given $p$ and $\eps$, there is  a constant $c$ such that if one camp has at least $n/2 + c$ individuals at the initial state, then
			it wins with probability at least $1- \eps$.  The surprising fact here is that $c$ does not depend on $n$, the population of the community. When $p=1/2$ and $\eps =.1$, one can set
			$c$ as small as 5.
			
			
			\begin{theorem} [The power of few]  \label{thm:const-advantage-ex}
				Consider the  (majority) process on $G(n, 1/2)$. Assume that the Red camp has at least  $n/2 + 5$  vertices at the initial state, where $n \ge 1000$. Then  Red wins after the fourth day with probability at least  $90\%$.
			\end{theorem}
			
			This result can be stated without the
			%
			\err\  %
			%
			$G(n, 1/2)$ %
			model; one can state an equivalent theorem by choosing the network uniformly, from the set of all graphs on $n$ vertices.
			
			This  result reveals an interesting phenomenon, which we call 
			``the power of few.''
			The collective outcome can be extremely sensitive, as a  modification of the smallest scale  in the initial setting leads to
			the opposite outcome.
			
			Our result applies in the following equivalent settings.
			
			\emph{Model 1.} We fix the two camps, of %
			%
			%
			sizes  %
			$n/2 +c$ and $n/2-c$,
			%
			and draw a random graph $G\sim G(n, p)$ over their union.
			
			\emph{Models 2.} We draw a random graph $G \sim G(n, p)$ first, let Red be a random subset of $n/2 +c$ vertices (chosen uniformly from all subsets of that size), and Blue be the rest.
			
			\emph{Model 3.} We split the society into two camps of size $n/2$ each, then draw the random graph $G\sim G(n, p)$ on their union, then recolor $c$
			%
			randomly selected  %
			Blue vertices to Red.
			
			%
			%
			\emph{Model 4.} Split the society into two camps (Red and Blue) of size $n/2-c$ each and a ``swing'' group (with no color yet) of $2c$ individuals. Draw the random graph on their union. Now let the swing group join the Red camp.
			
			With Model 3, we can imagine a balanced election process at the beginning. Then $c = 5$ people move to the other camp. \expref{Theorem}{thm:const-advantage-ex} asserts that this tiny group already %
			guarantees  %
			the final win with %
			overwhelming odds. Similarly, Model 4 implies that a swing group of size 10 decides the outcome.
			
			Our result can also be used to model the phenomenon that outcomes  in seemingly identical
			situations become opposite. Consider two communities, each has exactly $n$ individuals, sharing the same social network. In the first community,
			the  %
			Red camp  has
			size $n/2 +c$, and
			the  %
			Blue camp has $n/2-c$. In the second community,
			the  %
			Blue camp has $n/2+c$ and
			the  %
			Red camp has $n/2-c$. If $n$ is large,
			there is no way to tell the difference between the two communities.
			Even if we record everyone's initial opinion,  clerical errors will surely swallow the tiny difference
			of $2c$.  However, at the end, the collective opinion  will be opposite, with high probability.
			
			
			Now we state the general result for arbitrary constant density $p$.
			
			\begin{theorem}[Asymptotic bound] \label{thm:unbounded-advantage}
				Let $p$ be a constant in $(0, 1)$ and $c_n$ be a positive integer which may depend on $n$
				%
				.
				Assume that Red has $n/2 +c_n$ in day $0$ and the random graph is $G(n,p)$.
				Then Red wins after four days with probability at least $1 - K(p) \max \{n^{-1}, c_n^{-2} \} $, where $K(p)$ depends only on $p$.
			\end{theorem}
			
			Both  results follow from \expref{Theorem}{thm:const-advantage},  which, in a slightly technical form, describes how the process evolves day by day.
			Our results can be extended to cover the case when there are more than 2 opinions; details will appear in a later paper \cite{TranVusparse}.
			
			
			\subsection{Related results}
			
			Our problem is related to a %
			well-studied  %
			class of \emph{opinion exchange dynamics} problems. In the field of Computer Science, loosely-related processes are studied in \emph{population protocols} \cite{alistarh2015, alistarh2017}, where individuals/agents/nodes choose their next state based on that of their neighbors. The most separating difference is the network, as connections in these models often randomly change with time, while our study concerns a fixed network (randomly generated before the process begins).
			
			The survey by Mossel and Tamuz \cite{mossel2017} discussed %
			several  %
			models for these problems, including the \emph{DeGroot model} \cite{degroot1974}, where an individual's next state is a weighted average of its neighbors' current states, the \emph{voter model} \cite{clifford1973}, where individuals change states by emulating a random neighbor each day. The \emph{majority dynamics} model is in fact the same as ours, and is also more popular than the other two, having been studied in \cite{mossel2014, fountoulakis, benjamini}. The key difference, as compared to our study,  is in the %
			setups.  %
			In these earlier papers, each individual chooses his/her initial color uniformly at random. The central limit theorem thus guarantees that with high probability, the initial difference between the two camps is of order $\Theta (\sqrt n) $. Therefore, these papers did not touch upon the ``power of few'' phenomenon, which is our key message. On the other hand,  they considered sparse random graphs where the density $p = p_n\to 0$ as $n \to +\infty$.
			
			In \cite{benjamini}, Benjamini, Chan, O'Donnell, Tamuz, and Tan considered random graphs with $p \ge \lambda n^{-1/2} $, where $\lambda $ is a sufficiently large constant, and showed that the dominating color wins with probability at least $.4$ \cite[Theorem 1.2]{benjamini}, while conjecturing that this probability in fact tends to $1$ as $n \to \infty$. This conjecture was proved by Fountoulakis, Kang, and Makai
			\cite[Theorem 1.1]{fountoulakis}.
			
			\begin{theorem}[Fountoulakis, Kang, Makai]  \label{FKM}
				For any $0< \eps \le 1$ there is $\lambda = \lambda(\eps)$ such that the following holds for $p \ge \lambda n^{-1/2} $:
				With probability at least $1 -\eps$, over the choice of the random graph $G(n,p)$
				and the choice of the initial state, the dominating color wins after four days.
			\end{theorem}
			
			For related results on random regular graphs, see \cite{mossel2014, mossel2017}.
			
			
			\subsection{Extension for sparse random graphs}
			
			Note that the results presented in this paper only %
			apply to  %
			a constant $p$, which, in the context of $G(n, p)$, produces \emph{dense graphs}. For \emph{sparse graphs}, \ie, when $p = p_n$ depends on $n$ and tends to $0$ as $n \to +\infty$, the main ideas in this paper can be used, but with different algebraic techniques, to obtain a similar result.
			
			\begin{theorem} \label{thm:general-sparse}
				For any $0< \eps \le 1$ and $\lambda > 0$ there is $c = c (\eps, \lambda)$ such that for $p \ge (1 + \lambda) (\log n)/n$,
				if Red starts with $n/2 + c/p$ members, then it wins with probability at least $1 -\eps$.
			\end{theorem}
			
			The technical changes needed to prove this theorem require rewriting entire proofs with new computations, so we leave the proof to our future paper \cite{TranVusparse}. Additional information such as the length of the process and the explicit relation between the bound with $p$ and $c$ will also be discussed there.
			Notice that when $p$ is a constant, this result covers the ``Power of Few'' phenomenon as a special case, albeit with $c$ much larger than $5$. Thus, the techniques and results in this paper still have merit since they achieve a specific, surprisingly small constant.
			\expref{Theorem}{thm:general-sparse} no longer holds for $p < (\log n)/n$ as in this case there are, with high probability, isolated vertices. Any of these vertices  keeps %
			its  %
			original color forever. In this case, the number of Blue vertices  converges with time, and we obtain a bound on the limit in \cite{TranVusparse}.
			
			One can use  \expref{Theorem}{thm:general-sparse} to derive a ``delayed'' version  (in which Red may need more than 4 days) of \expref{Theorem}{FKM}, by first proving that with high probability, one side gains an advantage of size at least $C\sqrt{n}$ after the first day, for some constant $C$. This ``majority side'' then wins with high probability given $p \ge \lambda n^{-1/2}$ (which satisfies the requirement  %
			$p \ge (1 + \Omega(1))(\log n)/n$%
			) with $\lambda$ sufficiently large so that $\lambda C = pC\sqrt{n}$ is large. The detailed argument is in \expref{Appendix}{appx:FKM-proof}.
			
			
			\subsection {Notation}
			
			Below is the list of notations used throughout the paper.
			
			\begin{itemize}
				\item $\one_S$: membership function of the set $S\subset V$, so that $\one_S(v) \defeq \one_{v\in S}$.
				\item $R_t, B_t$: Respectively the sets of Red and Blue vertices after day $t$.
				(At this point everyone has updated their color $t$ times.)
				\item $\I_t(u) \ \defeq \ \bbm{1}_{\{u\in R_t\}}$: indicator of the event that $u$ is Red after day $t$, taking values in $\{0, 1\}$.
				\item $\J_t(u) \ \defeq \ 2\I_t(u) - 1$: indicator of the same event, but taking values in $\{-1, 1\}$.
				\item $u\adj v \equiv (u, v)\in E$: Event that $u$ and $v$ are adjacent.
				\item $\Neigh(v) \ \defeq \ \{u: u\adj v\}$: The neighborhood of $v$.
				\item $W_{uv} \ \defeq \ \mathbf{1}_{\{u\adj v\}}$ - Indicator of the adjacency between $u$ and $v$. %
				\item $\Phi(a, b) \defeq \int_a^b \frac{\eee^{-x^2/2}}{\sqrt{2\pi}}dx~$ and $\Phi(a) \defeq \Phi(-\infty, a), \ \Phi_0(a) \defeq \Phi(0, a)$.
				Note that $\Phi(a) = \Phi_0(a) + 1/2$.
				\item Given a random variable $X$, we let $\text{supp}(X)$ be the support of $X$, namely the smallest subset $S$ of $\mathbb{R}$ such that $\Pr(X\notin S) = 0$.
				When $X$ is discrete, $\text{supp}(X)$ is simply the set of possible values of $X$.
			\end{itemize}
			
			
			\subsection{Main Theorem} \label{subsec:main-theorem}
			
			The main theorem concerns dense graphs, where $p$ is at least a constant.  When given appropriate values for the parameters, it implies the ``Power of Few'' phenomenon in \expref{Theorem}{thm:const-advantage-ex}.
			
			\begin{theorem} \label{thm:const-advantage}
				Assume $p$ is a real number in $(0, 1)$ and $n$ and $c$ are positive integers such that 
				\begin{equation} \label{eq:n-p-c-conds}
					F(n, p, c) \defeq \sqrt{n} \Phi_0 \left( \frac{2pc %
						+ \min\{p, 1 - p\}  %
					}{\sqrt{p(1 - p)n}} \right) - .6\frac{1 - 2p + 2p^2}{\sqrt{p(1 - p)}} - \frac{.8}{p} > 0.
				\end{equation}
				Consider the election process on $G\sim G(n, p)$ with $\card{B_0}\le n/2 - c$.
				We have
				\begin{itemize}
					\item $\displaystyle \Pr\left(\card{B_1}\le \frac{n}{2} - \frac{.8\sqrt{n}}{p}\right) \ge 1 - P_1$, where
					\[
					P_1 = P_1(n, p, c) \ \defeq \ \frac{.25 + .7\left(1 - 2p + 2p^2\right)^2\left(1 + 4c^2n^{-2}\right)}{F^2(p, n, c)}.
					\]
					\item $\displaystyle \Pr\left(\card{B_2} \le .4n ~\Big|~ \card{B_1}\le \frac{n}{2} - \frac{.8\sqrt{n}}{p}\right) \ge 1 - P_2$, where
					\[
					P_2 = P_2(n) \ \defeq \ n^{-1}\exp\left( -(.076n - 1.38\sqrt{n}) \right).
					\]
					\item $\displaystyle \Pr\left(\card{B_3}\le \frac{p(n - 1)}{3} ~\Big|~ \card{B_2} \le .4n\right) \ge 1 - P_3$, where
					\[
					P_3 = P_3(n, p) \ \defeq \ n^{-1}\exp\left( -.02n(p^3n - 80) \right).
					\]
					\item $\displaystyle \Pr\left(\card{B_4} = 0 ~\Big|~ \card{B_3}\le \frac{p(n - 1)}{3}\right) \ge 1 - P_4$, where
					\[
					P_4 = P_4(n, p) \ \defeq \ n\exp\left(-\frac{2}{9}p^2(n-1)\right).
					\]
				\end{itemize}
				Consequently, $R_4 = V(G)$ with probability at least $1 - (P_1 + P_2 + P_3 + P_4)$.
			\end{theorem}
			
			%
			%
			%
			%
			Note that the condition \eqref{eq:n-p-c-conds} is asymptotically the same as $c = \Omega(p^{-3/2})$.
			
			The proof for this theorem has two main parts corresponding to the next two sections.
			\begin{enumerate}
				\item One day $1$, the number of Red and Blue neighbors of each node are binomial with means roughly 
				$n/2 + c$ and $n/2 - c$ respectively.
				The central limit theorem then implies that most of their masses are concentrated within an interval of length
				$\Theta (\sqrt n)$ around their respective expectations.
				A subinterval of constant length in this interval has $\Theta (n^{-1/2})$ mass.
				We thus expect that the probability that Reds outnumber Blues in 
				%
				the neighborhood of a given node %
				is %
				%
				$1/2 + \Omega (n^{-1/2})$. Thus, we expect that the number of Red nodes after the first day is $n/2 + \Omega (n^{1/2})$, meaning $n/2 - \Omega (n^{1/2})$ Blue nodes.\footnote[1]{While this paper was under review, Sah and Sawhney in a recent paper \cite{sahsawhney2021} showed that $\card{B_1}$ obeys a Gaussian limit law with mean $n/2 - \Theta(c\sqrt{pn})$ and variance $\Theta(n)$ when $n\to \infty$ for $p > (\log n)^{-1/16}$\,.}
				We use this intuition to find the right bound for $\card{B_1}$ in \expref{Theorem}{thm:const-advantage}, and prove it in \expref{Section}{sec:firstday-const}.
				
				\item On subsequent days, the coloring depends on the graph so the argument for day 1 is no longer available.
				We instead show that the number of Blues decreases %
				regardless of the coloring after day 1, as long as Red has $\Omega(\sqrt{n})$ extra nodes.
				Using Hoeffding bound (\expref{Theorem}{thm:hoeffding}), we show that for a given coloring, the probability one can find ``many'' nodes with a majority of neighbors in Blue is small, %
				sufficiently  %
				to beat the union bound over all colorings.
				This bypasses the dependency issue after day 1.
				The details are in \expref{Section}{sec:universal-win}.
			\end{enumerate}
			
			From \expref{Theorem}{thm:const-advantage}, one can deduce \expref{Theorems}{thm:const-advantage-ex} \expref{and}{thm:unbounded-advantage} in a few steps.
			%
			
			\begin{proof}[Proof of \expref{Theorem}{thm:const-advantage-ex}]
				Assume \expref{Theorem}{thm:const-advantage}. Observe that if  %
				\expeqref{Equation}{eq:n-p-c-conds} holds  %
				for some value of $n$, then  %
				it holds  %
				for all larger values of $n$.
				Let $n \ge 1000$, $p = 1/2$ and $c = 5$, we have  %
				\eqref{eq:n-p-c-conds}  %
				satisfied. A routine calculation then shows that
				\begin{linenomath}\begin{equation*}
						P_1(n, p, c)\le   %
						.0902  %
						, \quad \max\Bigl\{P_2(n), P_3(n, p), P_4(n, p)\Bigr\} < 10^{-10} , \end{equation*}\end{linenomath}
				\noindent which implies that $ \Pr \left(B_4\neq \varnothing \right) < .1$ or equivalently that Red wins in the fourth day with probability at least $.9$
				(conditioned on the event $\card{B_0} \le \frac{n}{2} - c$).
			\end{proof}
			
			\begin{proof}[Proof of \expref{Theorem}{thm:unbounded-advantage}]
				In this proof, only $n$ and $c =c_n$ can vary. We can assume, without loss of generality, that $c_n \le n/2$.
				Assuming \expref{Theorem}{thm:const-advantage},
				a routine calculation shows that $P_2, P_3, P_4 = \o(n^{-2})$ and, so $P_1 + P_2 + P_3 + P_4 = P_1 + \o(n^{-2})$.
				By \expref{Theorem}{thm:const-advantage}, we then have
				\[
				\Pr(R_4 = V(G)) \ge 1 - (P_1 + o(n^{-2}))
				= 1 - \frac{.25 + .7(1 - 2p + 2p^2)^2 \left(1 + 4c^2n^{-2}\right)}{\left(\sqrt{n}\ \Phi_0 \left( \frac{2pc_n  %
						+ \min\{p, 1 - p\}  %
					}{\sqrt{p(1 - p)n}} \right) - \frac{.6(1 - 2p + 2p^2)}{\sqrt{p(1 - p)}}- \frac{.8}{p}\right)^2} - o(n^{-2}).
				\]
				
				We make use of the fact $\Phi_0$ is increasing, while $x\mapsto \Phi_0(x)/x$ is decreasing,
				which implies $\Phi_0(x) \ge \min\{1, x/y\}\Phi_0(y) \ \forall x, y > 0$.
				Then:
				\begin{linenomath}\begin{equation*}
						\begin{aligned}
							\sqrt{n}\Phi_0\left(\frac{2pc_n  %
								+ \min\{p, 1 - p\}  %
							}{\sqrt{p(1 - p)n}}\right)
							\ge \sqrt{n}\min\left\{1, \frac{c_n}{\sqrt{(1 - p)n}} \right\} \Phi_0\left(2\sqrt{p}\right) \\
							\quad \ge \min\left\{\sqrt{n}, \frac{c_n}{\sqrt{1 - p}} \right\} \Phi_0\left(2\sqrt{p}\right)
							\ge \min\left\{\sqrt{n}, c_n \right\} \Phi_0\left(2\sqrt{p}\right).
						\end{aligned}
				\end{equation*}\end{linenomath}
				When $n$ and $c_n$ are both sufficiently large and $p$ is a constant, we have
				\begin{linenomath}\begin{equation*}
						\min\left\{\sqrt{n}, c_n \right\} \Phi_0\left(2\sqrt{p}\right) \ge 2\left[\frac{.6(1 - 2p + 2p^2)}{\sqrt{p(1 - p)}} + \frac{.8}{p} \right]
				\end{equation*}\end{linenomath}
				Therefore we can %
				give a lower bound on %
				%
				the denominator of $P_1$:
				\begin{linenomath}\begin{equation*}
						\sqrt{n}\ \Phi_0 \left( \frac{2pc_n  %
							+ \min\{p, 1 - p\}  %
						}{\sqrt{p(1 - p)n}} \right) - \frac{.6(1 - 2p + 2p^2)}{\sqrt{p(1 - p)}}- \frac{.8}{p}
						\ge \frac{\Phi_0\left(2\sqrt{p}\right)}{2}\min\left\{\sqrt{n}, c_n \right\}.
				\end{equation*}\end{linenomath}
				Now we use the trivial facts $c_n \le n/2$ and $1 - 2p + 2p^2 \le 1$ 
				%
				to give an upper bound on %
				the numerator of $P_1$:
				\begin{linenomath}\begin{equation*}
						.25 + .7(1 - 2p + 2p^2)^2 \left(1 + 4c^2n^{-2}\right) \le .25 + .7 \cdot 1 \cdot (1 + 1) < 2.
				\end{equation*}\end{linenomath}
				Combining the above, we get
				\begin{linenomath}\begin{equation*}
						P_1 \le \frac{2}{\left(\min\{c_n, \sqrt{n}\}\Phi_0(2\sqrt{p})/2 \right)^2}
						= \frac{4\Phi_0(2\sqrt{p})^{-2}}{\min\{\sqrt{n}, c_n\}^2}
						= 4\Phi_0(2\sqrt{p})^{-2} \max\left\{\frac{1}{n}, \frac{1}{c_n^2}\right\}.
				\end{equation*}\end{linenomath}
				The term $o(n^{-2})$ is also absorbed into this bound, since $n^{-2} < n^{-1}$.
				Therefore we have
				\begin{linenomath}\begin{equation*}
						\Pr(R_4 = V(G)) \ge 1 - K(p)\max\left\{\frac{1}{n}, \frac{1}{c_n^2}\right\}.
				\end{equation*}\end{linenomath}
				which is the desired bound where $K(p)$ can be chosen as $4\Phi_0(2\sqrt{p})^{-2} + 1$.
			\end{proof}
			
			
			%
			%
			%
			
			
			
			\subsection {Open questions}
			
			
			Let $\rho(p, k, n)$ be the probability that Red %
			wins  %
			if its camp has size $n/2 + k$ in the beginning.
			\expref{Theorem}{thm:const-advantage-ex} shows that $\rho (.5, 5, n) \ge .9$ (given that $n$ is sufficiently large).
			In other words, five defectors guarantee Red's victory with overwhelming odds when $p = 1/2$.
			In fact, we have $\rho(.5, 4, n) \ge .77$ by plugging in the same values for $\eps_1$ and $\eps_2$ with $c = 4$ in 
			%
			the proof of \expref{Theorem}{thm:const-advantage-ex}.  %
			We conjecture that one defector already brings a non-trivial advantage.
			
			\begin{conjecture}[The power of one] \label{conj:advantage-1}
				There is a constant $\delta(p) > 0$ for each constant $p > 0$ such that  $\rho (p, 1, n) \ge 1/2 + \delta(p)$ for all sufficiently large $n$.
				\footnote[1]{Both \expref{Conjectures}{conj:advantage-1} \expref{and}{conj:defectors} are proven for $p > (\log n)^{-1/16}$ in \cite{sahsawhney2021}. See footnote 1 a few paragraphs below \expref{Theorem}{thm:const-advantage}.}
			\end{conjecture}
			
			In the following numerical experiment,  we run $T = 10000$ independent trials.  In each trial, we  fix a set of $N = 10000$ nodes with $5001$ Red and $4999$  Blue (meaning $c=1$), generate a graph from $G(N, 1/2)$,  and simulate the process on the resulting graph. We record the number of wins  and the number of days to achieve the win in percentage  in \expref{Table}{tab:voting-sim}. Among others, we see that Red wins within $3$ days with frequency more than $.9$.
			The source code for the simulation along with execution
			instructions can be found online at \url{https://github.com/thbl2012/majority-dynamics-simulation}.
			%
			\begin{table}[h]
				\centering
				%
					%
					%
					%
					%
					%
					%
					%
					%
					%
					%
				
				\begin{tabular}{|ccccccrr|}\hline
					T & p & Red & Blue & 	Winner & Last day & Count & Frequency \\
					\hline
					$10^4$ & $1/2$ & $5001$ & $4999$ &	Blue & 3 & 496 & 4.96 \ \% \\
					$10^4$ & $1/2$ & $5001$ & $4999$ &	Blue & 4 & 77 & 0.77 \ \% \\
					$10^4$ & $1/2$ & $5001$ & $4999$ &	Blue & 5 & 3 & 0.03 \ \% \\
					$10^4$ & $1/2$ & $5001$ & $4999$ &	Blue & 7 & 1 & 0.01 \ \% \\
					$10^4$ & $1/2$ & $5001$ & $4999$ &	Red & 2 & 25 & 0.25 \ \% \\
					$10^4$ & $1/2$ & $5001$ & $4999$ &	Red & 3 & 9313 & 93.13 \ \% \\
					$10^4$ & $1/2$ & $5001$ & $4999$ &	Red & 4 & 85 & 0.85 \ \% \\
					\hline
				\end{tabular}
				\vspace{3pt}
				
				\caption{Winners and winning days with their frequencies}
				\label{tab:voting-sim}
			\end{table}
			
			Imagine that people defect from
			the  %
			Blue camp to
			the  %
			Red camp one by one. The \emph{value} of the
			%
			$i$-th  %
			defector is defined as $v(p, i,n) \defeq \rho (p, i,n) - \rho (p, i-1, n)$ (where we take $\rho (p, 0, n) = 1/2 $). It is intuitive to think that the value of each extra defector decreases. (Clearly defector number $n/2$ adds no value.)
			
			\begin{conjecture}[Values of defectors] \label{conj:defectors}
				For any fixed $p$, $i$, and $n$ large enough, $v(p, i,n) \ge v(p, i+1, n)$.\footnote[2]{See footnote 1 after \expref{Conjecture}{conj:advantage-1}.}
			\end{conjecture}
			
			
			It is clear that the \expref{Conjecture}{conj:defectors} implies \expref{Conjecture}{conj:advantage-1}, with $\delta = \frac{.4}{5} =.08 $, although the simulation results above suggests that $\delta$ can be at least $.43$.
			
			In the next two sections we give the proof of our main result, \expref{Theorem}{thm:const-advantage}.
			
			%
			%
			%
			%
			%
			%
			
			\section{Day One} \label{sec:firstday-const}
			
			
			Firstly, let us recall a few terms defined in \expref{Theorem}{thm:const-advantage}.
			\begin{linenomath}\begin{equation*}
					\begin{aligned}
						F(n, p, c) & \ \defeq \ \sqrt{n}\Phi_0\left(\frac{2pc  %
							+ \min\{p, 1 - p\}  %
						}{\sqrt{p(1 - p)n}} \right) - .6\frac{1 - 2p + 2p^2}{\sqrt{p(1 - p)}} - \frac{.8}{p}, \\
						P_1(n, p, c) & \ \defeq \ \frac{.25 + .7\left(1 - 2p + 2p^2\right)^2 \left(1 + 4c^2n^{-2}\right)}
						{F(n, p, c)^2}.
					\end{aligned}
			\end{equation*}\end{linenomath}
			
			\begin{lemma} \label{lem:const-advantage-day-1}
				Given $p \in (0, 1)$ and $n, c\in \bb{N}$ such that 
				$F(n, p, c)  > 0$,
				The election process on $G\sim G(n, p)$ with $\card{B_0}\le n/2 - c$ satisfies
				\[ \label{eq:const-advantage-day-1}
				\Pr \left(\card{B_1} > \frac{n}{2} - .8\frac{\sqrt{n}}{p}\right) \le P_1\left(n, p, c\right).
				\]
			\end{lemma}
			
			The core of the proof relies on some preliminary results regarding the difference of two binomial random variables, which we discuss next.
			
			
			\subsection{Background on  %
				the  %
				difference of Binomial Random Variables}
			
			The difference of two binomial random variables with the same probability $p$ can be written as a sum of independent random variables, each of which is either a $\Bin(1, p)$ variable or minus of one. A natural way to bound this sum is done via a Berry--Esseen normal approximation.
			
			\begin{theorem}[Berry--Esseen] \label{thm:esseen}
				Let $n$ be any positive integer. If $X_1, X_2, X_3, \ldots, X_n$ are random variables with zero means, variances $\sigma_1^2, \sigma_2^2, \cdots, \sigma_n^2 > 0$, and absolute third moments $\E{|X_i|^3} = \rho_i < \infty$, we have:
				\begin{linenomath}\begin{equation*} \label{BE1}
						\sup_{x\in \bb{R}}\abs{\Pr\left( \sum_{i=1}^nX_i \le x \right) - \Phi\left(\frac{x}{\sigma_X}\right)} \le C_0\cdot\frac{\sum_{i=1}^n\rho_i}{\sigma_X^3},
				\end{equation*}\end{linenomath}
				where $\sigma_X = \left(\sum_{i=1}^n\sigma_i^2\right)^{1/2}$ and $C_0$ is a constant.
			\end{theorem}
			
			\noindent The original proof by Esseen \cite{esseen} yielded  $C_0 = 7.59$, and this constant has been improved a number of times. The latest work by Shevtsova \cite{shevtsova} achieved  $C_0 = .56$, which will be used for the rest of the paper. A direct application of this theorem gives the following lemma.
			
			\begin{lemma} \label{thm:diff of binom r.v. gen.}
				For $p\in (0, 1)$, $\sigma = \sqrt{p(1-p)}$ and $n_1, n_2 \in \bb{N}$ such that $n_1 > n_2$.
				let $Y_1 \sim \Bin(n_1, p), \ Y_2 \sim \Bin(n_2, p)$ be independent random variables.
				Then for any $d\in \R$,
				\begin{linenomath}\begin{equation*}
						\Pr\left(Y_1 > Y_2 + d\right) \ge \frac{1}{2} + \Phi_0\left(\frac{p(n_1 - n_2) - d}{\sqrt{p(1 - p)(n_1 + n_2)}}\right) - \frac{.56\left(1 - 2p + 2p^2\right)}{\sqrt{p(1 - p)(n_1 + n_2)}}.
				\end{equation*}\end{linenomath}
			\end{lemma}
			
			\begin{proof}
				Let $X = Y_1 - Y_2$. By definition, we have
				$
				X = X_1 + X_2 + X_3 + \cdots + X_{n_1+n_2},
				$
				where all %
				the $X_i$   %
				are independent and either $X_i$ or $-X_i$ is $\Bin(1, p)$.
				Then $\E{X} = \sum_i\E{X_i} = p(n_1 - n_2)$.
				For all $i$,
				\begin{linenomath}\begin{equation*}
						\Var{X_i} = p(1 - p) \ \text{ and } \ \E{|X_i - \E{X_i}|^3} = p(1 - p)^3 + (1 - p)p^3 = p(1 - p)(1 - 2p + 2p^2).
				\end{equation*}\end{linenomath}
				Applying \expref{Theorem}{thm:esseen}, we have
				\begin{linenomath}\begin{equation*} \label{eq:temp68}
						\begin{aligned}
							& \Pr\left(Y_1 \le Y_2 + d\right) \ = \ \Pr\left(X - \E{X} \le d - p(n_1 - n_2)\right)
							\\
							& \ \le \ \Phi\left(\frac{d - p(n_1 - n_2)}{\sqrt{\Var{X}}} \right)
							+ .56 \frac{\sum_i \E{|X_i - \E{X_i}|^3} }{\Var{X}^{3/2}}
							\\
							& \ = \ \Phi\left(\frac{d - p(n_1 - n_2)}{\sqrt{p(1 - p)(n_1 + n_2)}} \right)
							+ .56 \frac{p(1 - p) (1 - 2p + 2p^2)(n_1 + n_2) }{(p(1 - p)(n_1 + n_2))^{3/2}}
							\\
							& \ = \ \frac{1}{2} - \Phi_0\left(\frac{p(n_1 - n_2) - d}{\sqrt{p(1 - p)(n_1 + n_2)}} \right)
							+ \frac{.56(1 - 2p + 2p^2) }{\sqrt{p(1 - p)(n_1 + n_2)}},
						\end{aligned}
				\end{equation*}\end{linenomath}
				and the claim follows by taking the complement event.
			\end{proof}
			
			%
			\begin{lemma} \label{thm:1/sqrt{n}}
				Let $p \in (0, 1)$ be a constant, and $X_1 \sim \Bin(n_1, p)$ and $X_2 \sim \Bin(n_2, p)$ be independent r.v.s. Then for any integer $d$,
				\begin{linenomath}\begin{equation*}
						\Pr\left(X_1 = X_2 + d \right) \le \frac{1.12\left(1 - 2p + 2p^2\right)}{\sqrt{p(1 - p)(n_1 + n_2)}}.
				\end{equation*}\end{linenomath}
			\end{lemma}
			
			\begin{proof}
				Let $n = n_1+n_2$ and $\mu = \E{X_1} - \E{X_2} = p(n_1 - n_2)$. Fix $\varepsilon \in (0, 1)$, by the same computations in \expref{Lemma}{thm:diff of binom r.v. gen.}, we have
				\begin{linenomath}\begin{equation*} \label{eq:temp001}
						\begin{aligned}
							& \Pr\left(X_1 - X_2 \le d - \varepsilon\right) \ge
							\Phi\left(\frac{d - \mu - \varepsilon}{\sqrt{p(1 - p)n}}\right)
							- \frac{.56\left(1 - 2p + 2p^2\right)}{\sqrt{p(1 - p)n}}, \\
							%
					%
					%
							& \Pr\left(X_1 - X_2 < d + \varepsilon\right) \le
							\Phi\left(\frac{d - \mu + \varepsilon}{\sqrt{p(1 - p)n}}\right)
							+ \frac{.56\left(1 - 2p - 2p^2\right)}{\sqrt{p(1 - p)n}}.
						\end{aligned}
				\end{equation*}\end{linenomath}
				It follows that
				\begin{linenomath}\begin{equation*}
						\begin{aligned}
							& \Pr(X_1 = X_2 + d)
							\ \le \ \Pr(d - \varepsilon < X_1 - X_2 < d + \varepsilon)\\
							& \ \le \ \Phi\left(\frac{d - \mu + \varepsilon}{\sqrt{p(1 - p)n}}\right) -
							\Phi\left(\frac{d - \mu - \varepsilon}{\sqrt{p(1 - p)n}}\right) + \frac{1.12\left(1 - 2p + 2p^2\right)}{\sqrt{p(1 - p)n}}.
						\end{aligned}
				\end{equation*}\end{linenomath}
				Letting $\varepsilon \to 0$, we obtain the desired claim.
			\end{proof}
			
			\subsection*{Proof of \expref{Lemma}{lem:const-advantage-day-1}}
			
			Recall that $\card{R_1} = n - \card{B_1}$. Our goal is to 
			%
			give an upper bound on  %
			the probability that $\card{R_1} < n/2 + d\sqrt{n}$ for any given term $d$.
			Recall that $\card{R_1} = \sum_{v\in V}\I_1(v)$, where $\I_1(v) $ is $1$ if $v$ is Red after Day 1 and $0$ otherwise. We have:
			Since the indicators are not independent, a natural choice for bounding their sum is to use Chebysev's inequality. We proceed in two steps:
			\begin{enumerate}
				%
				\item Give a lower bound on $\E{\card{R_1}}$ by lowerbounding  %
				each term $\E{\I_1(v)}$.  
				%
				\item Give an upper bound on $\Var{\card{R_1}}$ by upperbounding %
				each $\Var{\I_1(v)}$ and $\Cov{\I_1(v), \I_1(v')}$.
			\end{enumerate}
			Let $\sigma \defeq \sqrt{p(1 - p)}$, which appears in most equations in this proof.
			Note that $1 - 2p + 2p^2 = 1 - 2\sigma^2$.
			
			%
			%
			
			%
			%
			
			\begin{claim} \label{claim:day-1-expection}
				$ \displaystyle
				\E{\card{R_1}} \ge \frac{n}{2} + F_1(n, p, c)\sqrt{n},
				$ where
				\[
				F_1(n, p, c) = \sqrt{n}\Phi_0\left(\frac{2pc  %
					+ \min\{p, 1 - p\}  %
				}{\sigma\sqrt{n}} \right) - .6\frac{1 - 2\sigma^2}{\sigma} = F(n, p, c) + \frac{.8\sqrt{n}}{p}.
				\]
				
			\end{claim}
			\begin{proof}
				%
				For $v\in V$ and $S\subset V$, let $d_S(v)$ be the number of neighbors $v$ has in $S$.
				%
				%
				By the majority and tie-breaking rules  %
				, %
				we have for each $v\in V$,
				\begin{equation} \label{eq:a_1 meaning}
					v \in R_1 \ \iff \ d_{R_0}(v) > d_{B_0}(v) - \I_0(v).
				\end{equation}
				Note that $d_{R_0}(v) \sim \Bin\left(\card{R_0} - \I_0(v), p\right)$ and $d_{B_0}(v) \sim \Bin\left(\card{B_0} + \I_0(v) - 1, p\right)$.
				By \expref{Lemma}{thm:diff of binom r.v. gen.}, we have:
				\begin{linenomath}\begin{equation*} \label{eq:i1-v-bound}
						\begin{aligned}
							& \E{\I_1(v)} = \Pr\left(v\in R_1\right) = \Pr\left(d_{R_0}(v) > d_{B_0}(v) - \I_0(v) \right) \\
							& \ge \frac{1}{2} + \Phi_0\left(\frac{p\left(\card{R_0} - \card{B_0} + 1 - 2\I_0(v)\right) + \I_0(v)}{\sigma\sqrt{n - 1}}\right)
							- \frac{.56\left(1 - 2\sigma^2\right)}{\sigma\sqrt{n - 1}} \\
							& = \frac{1}{2} + \Phi_0\left(\frac{2pc + p_v}{\sigma\sqrt{n - 1}}\right)
							- \frac{.56\left(1 - 2\sigma^2\right)}{\sigma\sqrt{n - 1}}
							%
							%
							\ge \frac{1}{2} + \Phi_0\left(\frac{2pc + \min\{p, 1 - p\}}{\sigma\sqrt{n}}\right)
							- \frac{.6\left(1 - 2\sigma^2\right)}{\sigma\sqrt{n1}},  %
						\end{aligned}
				\end{equation*}\end{linenomath}
				where
				$ \displaystyle
				p_v \defeq p\left(1 - \I_0(v)\right) + (1 - p)\I_0(v) \ge %
				\min\{p, 1 - p\}  %
				$, hence the last inequality.
				Summing over all $v\in V$, we get
				\begin{linenomath}\begin{equation*}
						\begin{aligned}
							\E{\card{R_1}}
							%
							%
							\ge \frac{n}{2} + \left[\sqrt{n}\Phi_0\left(\frac{2pc  %
								+ \min\{p, 1 - p\}  %
							}{\sigma\sqrt{n}}\right)
							- \frac{.6\left(1 - 2\sigma^2\right)}{\sqrt{p(1 - p)}} \right] \sqrt{n}.
						\end{aligned}
				\end{equation*}\end{linenomath}
				The proof is complete.
			\end{proof}
			
			%
			%
			
			\begin{claim} \label{claim:day-1-covariance}
				$ \displaystyle
				\Var{\card{R_1}} \ \le \ .25n + .7\left(1 - 2\sigma^2\right)^2 \left(n + 4c^2n^{-1}\right).
				$
			\end{claim}
			
			\begin{proof}
				We first have:
				\begin{linenomath}\begin{equation*}
						\Var{\card{R_1}} = \sum_{v\in V}\Var{\I_1(v)} + 2\sum_{v_1\neq v_2}\Cov{\I_1(v_1), \I_1(v_2)}.
				\end{equation*}\end{linenomath}
				\noindent The variances $\Var{\I_1(v)}$ are easy due to $\I_1(v)$ being a  Bernoulli variable:
				\begin{equation} \label{eq:day-1-variance}
					\Var{\I_1(v)} = \E{\I_1(v)}\bigl(1 - \E{\I_1(v)}\bigr) = .25 - \left(\E{\I_1(v)} - .5\right)^2
					\le .25.
				\end{equation}
				Bounding the covariance $\Cov{\I_1(v_1), \I_1(v_2)}$ for two distinct vertices $v_1, v_2$ requires a bit more care, as the indicators are dependent.
				%
				By definition
				\begin{linenomath}\begin{equation*}
						\Cov{\I_1(v_1), \I_1(v_2)} \ = \ \Pr\left(v_1, v_2 \in R_1\right) - \Pr\left(v_1 \in R_1\right) \Pr\left(v_2 \in R_1\right).
				\end{equation*}\end{linenomath}
				Consider the event $\{v_1, v_2 \in R_1\}$;  $\Pr\left(v_1, v_2 \in R_1\right)$ can be written as
				\begin{linenomath}\begin{equation*}
						\Pr\left(v_1, v_2 \in R_1 \given v_1 \adj v_2\right)\Pr(v_1 \adj v_2) + \ \Pr\left(v_1, v_2 \in R_1 \given v_1 \not\adj v_2\right)\Pr(v_1 \not\adj v_2).
				\end{equation*}\end{linenomath}
				Notice  that after we reveal the adjacency between $v_1$ and $v_2$, the remaining vertices in the neighborhoods of $v_1$ and $v_2$ are independent. Letting $a_i = \Pr(v_i\in R_1\mid v_1\adj v_2)$, $b_i = \Pr(v_i\in R_1\mid v_1\not\adj v_2)$ %
				, we have
				%
				%
				%
				%
				%
				%
				%
				\begin{linenomath}\begin{equation*}
						\Pr\left(v_1, v_2 \in R_1\right) =  pa_1a_2 + (1 - p)b_1b_2.
				\end{equation*}\end{linenomath}
				%
				%
				Splitting up the other two events similarly gives
				$\Pr\left(v_i \in R_1\right) =  pa_i + (1 - p)b_i$
				for $i = 1, 2$.
				Putting the above together, we can write $\Cov{\I_1(v_1), \I_1(v_2)}$ as
				\begin{equation} \label{eq:temp34}
					\begin{aligned}
						\ & \ pa_1a_2 + (1 - p)b_2b_2 - \left(pa_1 + (1 - p)b_1\right) \left(pa_2 + (1 - p)b_2 \right) \\
						\ & = \ p(1 - p)(a_1 - b_1)(a_2 - b_2) = \sigma^2(a_1 - b_1)(a_2 - b_2).
					\end{aligned}
				\end{equation}
				To bound $a_i - b_i$, we use logical reasoning to express it as something separate from the current context.
				We start with $i = 1$. The analysis for $a_2 - b_2$ is analogous.
				%
				Assume  %
				$v_1 \sim v_2$ for now, then
				\begin{linenomath}\begin{equation*}
						a_1 = \Pr\left(d_{R_0}(v_1) \ge d_{B_0}(v_1) - \I_0(v_1) + 1\right),
				\end{equation*}\end{linenomath}
				by Eq. \eqref{eq:a_1 meaning}.
				Define two binomial random variables:
				\begin{linenomath}\begin{equation*}
						\begin{aligned}
							X & = d_{R_0}(v_1) - \I_0(v_2) = \card{\Neigh(v_1)\cap  R_0\setminus \{v_1, v_2\}} 
							\sim \Bin\Bigl(\card{R_0} - \I_0(v_1) - \I_0(v_2), p\Bigr), \\
							Y & = d_{B_0}(v_1) + \I_0(v_2) - 1 = \card{\Neigh(v_1)\cap B_0\setminus \{v_1, v_2\}}
							\sim \Bin\Bigl(\card{B_0} + \I_0(v_1) + \I_0(v_2) - 2, p\Bigr).
						\end{aligned}
				\end{equation*}\end{linenomath}
				We have
				\begin{linenomath}\begin{equation*}
						\begin{aligned}
							a_1 = \Pr\left(X + \I_0(v_2) + \I_0(v_1) \ge Y + 2 - \I_0(v_2)\right)
							= \Pr\left(X - Y \ge 1 - \I_0(v_1) - \J_0(v_2)\right).
						\end{aligned}
				\end{equation*}\end{linenomath}
				Similarly, when conditioning on $v_1\not\adj v_2$, we get
				$
				b_1 = \Pr\left(X - Y \ge 1 - \I_0(v_1)\right).
				$
				
				\noindent Let us perform  %
				a %
				case analysis on 
				%
				the initial color of $v_2$.  %
				\begin{linenomath}\begin{equation*}
						\begin{aligned}
							v_2 \in R_0
							& ~\Rightarrow~
							a_1 - b_1 = \Pr(X - Y \ge - \I_0(v_1)) - \Pr(X - Y \ge 1 - \I_0(v_1)) = \Pr(X - Y = - \I_0(v_1)).
							\\
							v_2 \in B_0
							& ~\Rightarrow~
							a_1 - b_1 = \Pr(X - Y \ge 2 - \I_0(v_1)) - \Pr(X - Y \ge 1 - \I_0(v_1)) = -\Pr(X - Y = 1 - \I_0(v_1)).
						\end{aligned}
				\end{equation*}\end{linenomath}
				%
						%
							%
							%
							%
						%
				\noindent Both cases give us
				\begin{linenomath}\begin{equation*}
						a_1 - b_1 = \J_0(v_2)\Pr\Bigl(X - Y = 1 - \I_0(v_1) - \I_0(v_2)\Bigr),
				\end{equation*}\end{linenomath}
				where $X\sim \Bin\Bigl(\card{R_0} - \I_0(v_1) - \I_0(v_2), p\Bigr)$
				and $Y\sim \Bin\Bigl(\card{B_0} + \I_0(v_1) + \I_0(v_2) - 2, p\Bigr)$.
				
				\noindent We apply the same analysis for $a_2$ and $b_2$ and use Eq. \eqref{eq:temp34} to get
				\begin{linenomath}\begin{equation*}
						\Cov{\I_1(v_1), \I_1(v_2)} = \sigma^2 \J_0(v_1)\J_0(v_2) \Pr\Bigl(X - Y = 1 - \I_0(v_1) - \I_0(v_2)\Bigr)^2.
				\end{equation*}\end{linenomath}
				%
				If the initial colors of $v_1$ and $v_2$ %
				are %
				different  %
				$\Cov{\I_1(v_1), \I_1(v_2)} \le  0$,
				so we can exclude them from the upper bound.
				When $v_1$ and $v_2$ are of the same initial color, by \expref{Lemma}{thm:1/sqrt{n}}, we have
				\begin{linenomath}\begin{equation*}
						\Cov{\I_1(v_1), \I_1(v_2)} \le \sigma^2 \left(\frac{1.12\left(1 - 2\sigma^2\right)}{\sigma\sqrt{n - 2}}\right)^2
						%
						\le 1.4\frac{\left(1 - 2\sigma^2\right)^2}{n}.
				\end{equation*}\end{linenomath}
				Hence
				\begin{linenomath}\begin{equation*} \label{eq:day-1-covariance}
						\begin{aligned}
							& \sum_{v_1\neq v_2} \Cov{\I_1(v_1), \I_1(v_2)}
							\le \left[\binom{\card{R_0}}{2} + \binom{\card{B_0}}{2}\right]
							\cdot 1.4\cdot \frac{\left(1 - 2\sigma^2\right)^2}{n} \\
							& \ = \frac{1.4}{n} \left[\frac{n^2}{4} + c^2 - \frac{n}{2}\right] \left(1 - 2\sigma^2\right)^2
							\le .35\left(1 - 2\sigma^2\right)^2 \left(n + \frac{4c^2}{n}\right).
						\end{aligned}
				\end{equation*}\end{linenomath}
				\expeqref{Equations}{eq:day-1-variance} \expeqref{and}{eq:day-1-covariance} together yield
				\begin{equation} \label{eq:variance-bound}
					\Var{\card{R_1}} \le .25n + .7\left(1 - 2\sigma^2\right)^2 \left(n + 4c^2n^{-1}\right).
				\end{equation}
				The proof is complete.
			\end{proof}
			
			
			\noindent From \expref{Claims}{claim:day-1-expection} \expref{and}{claim:day-1-covariance}, %
			%
			Chebyshev's inequality gives
			%
			%
			%
			%
			%
			%
			%
			%
			\begin{linenomath}\begin{equation*} \label{eq:chebyshev-R1-bound}
					\begin{aligned}
						& \Pr \left(\card{B_1} > \frac{n}{2} - \frac{.8}{p}\sqrt{n} \right)
						\ = \ \Pr \left(\card{R_1} < \frac{n}{2} + \frac{.8}{p}\sqrt{n} \right) \\
						& \ \le \ \frac{\Var{\card{R_1}}}{\left(\E{\card{R_1}} - \frac{n}{2} - \frac{.8}{p}\sqrt{n}\right)^2}
						\ \le \ \frac{.25n + .7\left(1 - 2\sigma^2\right)^2 \left(n + 4c^2n^{-1}\right)}
						{n\left(\sqrt{n}\Phi_0\left(\frac{2pc  %
								+ \min\{p, 1 - p\}  %
							}{\sigma\sqrt{n}}\right) - .6\frac{1 - 2\sigma^2}{\sqrt{p(1 - p)}} - \frac{.8}{p} \right)^2}.
					\end{aligned}
			\end{equation*}\end{linenomath}
			Dividing both the nominator and denominator by $n$ and substituting $\sigma = \sqrt{p(1 - p)}$, we get back $P_1(n, p, c)$ and finish
			the proof of \expref{Lemma}{lem:const-advantage-day-1}.
			art one of \expref{Theorem}{thm:const-advantage} is complete.
			%
			%
			%
			%
			%
			
			%
			%
			%
			%
			%
			%
			\section{Day Two and after} \label{sec:universal-win}
			
			\noindent Next, we analyze the situation after the first day.
			%
			As mentioned in \expref{Section}{subsec:main-theorem}, the previous %
			argument  %
			is unusable due to the loss of independence from day $2$ onwards.
			Instead, we use ``shrinking arguments'' to argue that it is likely for the Blue camp to repeatedly shrink to empty, regardless of the choice of its members, due to %
			the structure of $G$.  %
			The core of our shrinking argument is \emph{Hoeffding's inequality}, a classical result that gives exponentially small probability tails for sums of independent random variables.
			\begin{theorem}[Hoeffding's inequality]  \label{thm:hoeffding}
				Let $\{X_i\}_{i=1}^n$ be independent random variables and $\{a_i\}_{i=1}^n$, $\{b_i\}_{i=1}^n$ be real numbers such that for all $i=\overline{1, n}$, $\text{supp}(X)\subseteq [a_i, b_i]$. Then for $X = X_1+X_2+\cdots + X_n$, we have
				\begin{linenomath}\begin{equation*}
						\max \Bigl\{\Pr\left( X \ge \E{X} + t \right),\ \Pr\left( X \le \E{X} -t \right)^{} \Bigr\} \le \exp\left(-\frac{2t^2}{\sum_{i=1}^n (b_i - a_i)^2}\right).
				\end{equation*}\end{linenomath}
			\end{theorem}
			The proof of Hoeffding's inequality is 
			available in %
			several  %
			graduate level probability textbooks, \eg, \cite{vershynin2018}. The original proof given by Hoeffding appeared in \cite{hoeffding1963}.
			
			Before discussing our shrinking arguments, we quickly record a useful
			%
			%
			%
			%
			lemma to simplify binomial coefficients that may appear in union bounds.
			%
			%
			%
			%
			%
			\begin{lemma}  \label{lem:binom-bound}
				For any integers $n\ge 1$ and $0\le k\le n$, $\displaystyle \binom{n}{k} < \frac{2^n}{\sqrt{n}}$.
			\end{lemma}
			%
			\begin{proof}
				This is well known.  Observing that $\binom{n}{k}\le \binom{n}{\lfloor n/2\rfloor}$,
				the case of even $n$ is stated in \cite[2.6.2 Proposition]{invitation}.  The case of
				odd $n=2m-1$ follows immediately:
				\begin{linenomath}\begin{equation*}
						\binom{2m - 1}{m - 1} = \frac{1}{2} \binom{2m}{m} < \frac{1}{2} \frac{2^{2m}}{\sqrt{2m}}
						= \frac{2^{2m - 1}}{\sqrt{2m}} < \frac{2^{2m - 1}}{\sqrt{2m - 1}}.  \qedhere
				\end{equation*}\end{linenomath}
			\end{proof}
			
			A simple yet useful shrinking argument is that, in the $G(n, p)$ model,  it is with high probability that all vertices in $G$ have many neighbors, so a small enough Blue camp will not be able to influence anyone by a majority, thus inevitably vanishes the next day.
			
			\begin{lemma}[Part 4 of \expref{Theorem}{thm:const-advantage}] \label{prop:uni-large-linear-advantage}
				For $p\in (0, 1)$ and $n\in \bb{N}_{>1}$, with probability at least
				\begin{linenomath}\begin{equation*}
						1 - n \exp\left(-\frac{2}{9}p^2(n-1)\right) = 1 - P_4(n, p),
				\end{equation*}\end{linenomath}
				$G$ is such that all vertices have more than $\frac{2}{3}p(n - 1)$ neighbors, thus any choice of the Blue camp of at most $\frac{1}{3}p(n - 1)$ members vanishes the next day.
			\end{lemma}
			
			\begin{proof}
				In a $G(n, p)$ graph, $d(v)$ is a sum of $(n-1)$ $\Bin(1, p)$ random variables, so \expref{Theorem}{thm:hoeffding} implies that for any $u\in V$,
				\begin{linenomath}\begin{equation*}
						\Pr\left(d(u)\le \frac{2}{3}p(n-1)\right)
						\le \Pr\left(d(u) - \E{d(u)} \le -\frac{1}{3}p(n-1)\right) \le \exp\left(-\frac{2}{9}p^2(n-1)\right).
				\end{equation*}\end{linenomath}
				By a union bound, the probability that all vertices have more than $\frac{2}{3}p(n-1)$ neighbors is at least $1 - n\exp\left(-\frac{2}{9}p^2(n-1)\right) = 1 - P_4(n, p)$. Given this, a Blue camp of size $\frac{1}{3}p(n-1)$ surely vanishes the next day since it cannot form a majority in 
				%
				the neighborhood of any vertex.  %
				The result then follows.
			\end{proof}
			
			This simple lemma completes the fourth part of \expref{Theorem}{thm:const-advantage}.
			The arguments for Parts 2 and 3 require slightly more complexity, but both come directly from the following lemma.
			
			
			
			\begin{lemma} \label{lem:uni-bad-set-bound}
				Let $p\in (0, 1)$, $n, n_0 \in \mathbb{N}$, $n_0 < \frac{n}{2}$. Then for all $m\in \mathbb{N}, m\le n$, with probability at least
				\begin{linenomath}\begin{equation*}
						1 - \frac{4^n}{n} \exp\left(-\frac{2p^2(n - 2n_0 - 1)^2m}{n + m - 2} \right),
				\end{equation*}\end{linenomath}
				$G$ is such that any choice of the Blue camp of at most $n_0$ members shrinks to below $m$ in the next day.
			\end{lemma}
			
			\begin{proof}[Proof of \expref{Lemma}{lem:uni-bad-set-bound}]
				Consider a subset $S$ of $V$ with $m$ elements.
				We first bound the probability that $S$ entirely turns Blue the next day.
				Let $(R, B)$ be the initial coloring with $\card{B} = n_0 < n/2$.
				Similarly to the usual $\J_t$, let $\J(v) = 1$ if $v\in R$ and $-1$ otherwise.
				For each $v\in V$, let
				$\dif(v) \defeq \card{\Neigh(v) \cap R} - \card{\Neigh(v) \cap B} = \sum_{u\in V} \J(u)W_{uv}$,
				and let $\dif(S) \defeq \sum_{v\in S}\dif(v)$.
				%
				%
					%
						%
						%
						%
					%
				%
				%
						%
						%
						%
						%
				%
				%
				We break down $\dif(S)$ as follows.
				\begin{equation} \label{eq:difS}
					\begin{aligned}
						\dif(S) = \sum_{v\in V} \Bigl[\one_S(v)\sum_{u\in V}\J(u)W_{uv}\Bigr]
						= \sum_{u \neq v} \Bigl[\J(u)\one_S(v) + \J(v)\one_S(u) \Bigr] W_{uv}.
					\end{aligned}
				\end{equation}
				This is now a sum of independent variables, so we can apply \expref{Theorem}{thm:hoeffding}.
				\begin{equation}  \label{eq:uni-bad-set-bound-temp1}
					\begin{aligned}
						\Pr\left(S\subseteq B_1\mid B_0\right)
						\ \le \ \Pr(\dif(S)\le 0)
						%
						\ \le \ \exp\left[
						- \frac{\E{\dif(S)}^2}
						{\sum_{u\neq v} \bigl(\sup \text{supp}(X_{uv}W_{uv})
							- \inf \text{supp}(X_{uv}W_{uv})\bigr)^2}
						\right],
					\end{aligned}
				\end{equation}
				where $X_{uv} \defeq \J(u)\one_S(v) + \J(v)\one_S(u)$.
				The following table sums up respective values of $X_{uv}$ for $(u, v)$ among the sets
				$R\cap S$, $R\setminus S$, $B\cap S$ and $B\setminus S$.
				\begin{equation}  \label{eq:X_(uv)-table}
					\begin{array}{c|c|c|c|c}
						u\setminus v & R\cap S               & B\cap S               & R\cap S^c & B\cap S^c \\ \hline
						R\cap S      & 2 & 0 & \multirow{2}{*}{$1$} & \multirow{2}{*}{$-1$} \\ \cline{1-3}
						B\cap S      & 0 & -2 &           &  \\ \hline
						R\cap S^c    & \multicolumn{2}{|c|}{1}               & \multicolumn{2}{|c}{\multirow{2}{*}{0}}  \\ \cline{1-3}
						B\cap S^c    & \multicolumn{2}{|c|}{-1}              & \multicolumn{2}{|c}{}
					\end{array}
				\end{equation}
				
				\noindent Substituting back into \expeqref{Equation}{eq:difS}, we get
				\begin{linenomath}\begin{equation*}
						\dif(S) \ = \sum_{\{u, v\} \subset S\cap R}\hspace{-1em}\left(2W_{uv}\right)
						- \sum_{\{u, v\} \subset S\cap B}\hspace{-1em}\left(2 W_{uv}\right)
						+ \sum_{u\in S} \sum_{v\in R\setminus S}W_{uv} - \sum_{u\in S} \sum_{v\in B\setminus S}W_{uv} .
				\end{equation*}\end{linenomath}
				For convenience in the following computations, let $m_R\defeq \card{S\cap R}$ and $m_B \defeq \card{S\cap B}$.
				
				\noindent Note that $\card{B\setminus S} = n_0 - m_B$ and $\card{R\setminus S} = n - n_0 - m_R$.
				Let $W$ be a $\Bin(1, p)$ random variable.
				Observe that $S$ is a sum of
				$\binom{m_R}{2} = \frac{m_R(m_R - 1)}{2}$ copies of $2W$,
				$\binom{m_B}{2} = \frac{m_B(m_B - 1)}{2}$ copies of $(-2W)$,
				$m(n - n_0 - m_R)$ copies of $W$,
				and $m(n_0 - m_B)$ copies of $(-W)$, all independent.
				Therefore
				\begin{linenomath}\begin{equation*}  \label{eq:temp97}
						\begin{aligned}
							\E{\dif(S)}
							& = p\left[2\cdot\frac{m_R(m_R - 1)}{2} - 2\cdot\frac{m_B(m_B - 1)}{2} + m(n - n_0 - m_R) - m(n_0 - m_B) \right] \\
							& = p\Bigl[m_R^2 - m_B^2 - m_R + m_B + m(n - 2n_0 - m_R + m_B)\Bigr] \\
							& = p\Bigl[(m_R - m_B)(m_R + m_B) - m(m_R - m_B) + m(n - 2n_0) - m_R + m_B)\Bigr] \\
							& = p\Bigl[m(m_R - m_B) - m(m_R - m_B) + m(n - 2n_0) - (m - m_B) + m_B\Bigr] \\
							& = p\Bigl[m(n - 2n_0) - m + 2m_B\Bigr] \ge p\Bigl[m(n - 2n_0) - m\Bigr] = pm(n - 2n_0 - 1).
						\end{aligned}
				\end{equation*}\end{linenomath}
				Moreover, $\text{supp}(W_{uv}) = {0, 1}$, so $\text{supp}(X_{uv}W_{uv}) = \{0, X_{uv}\}$, thus
				\[
				\sup \text{supp}(X_{uv}W_{uv}) - \inf \text{supp}(X_{uv}W_{uv}) = |X_{uv}| \ \text{ for all } u\neq v.
				\]
				%
				Therefore the denominator of the exponent in \expeqref{Equation}{eq:uni-bad-set-bound-temp1} is
				\begin{linenomath}\begin{equation*}
						\begin{aligned}
							& \sum_{u\neq v} \bigl(\sup \text{supp}(X_{uv}W_{uv})
							- \inf \text{supp}(X_{uv}W_{uv})\bigr)^2
							= \sum_{u\neq v} X_{uv}^2 \\
							& \ = \ \sum_{\{u, v\} \subset S\cap R}\hspace{-1em}2^2
							+ \sum_{\{u, v\} \subset S\cap B}\hspace{-1em}(-2)^2
							+ \sum_{u\in S} \sum_{v\in R\setminus S}1^2
							+ \sum_{u\in S} \sum_{v\in B\setminus S}(-1)^2 \\
							& \ = \ 4\cdot\frac{m_R(m_R - 1)}{2} + 4\cdot\frac{m_B(m_B - 1)}{2} + m(n - n_0 - m_R) + m(n_0 - m_B) \\
							& \ = \ 2m_R^2 + 2m_B^2 - 2m_R - 2m_B + m(n - m_R - m_B)
							= 2(m_R + m_B)^2 - 4m_Rm_B - 2m + m(n - m) \\
							& \ \le \ 2(m_R + m_B)^2 - 2m + m(n - m)
							= 2m^2 - 2m + m(n - m) = m(n + m - 2).
						\end{aligned}
				\end{equation*}\end{linenomath}
				Substituting this bound into \expeqref{Equation}{eq:uni-bad-set-bound-temp1}, we get
				\begin{linenomath}\begin{equation*}
						\Pr(S\subseteq B_1\mid B_0 = B)
						\le \exp\left[-\frac{2\bigl(pm(n - 2n_0 - 1)\bigr)^2}{m(n + m - 2)}\right]
						= \exp\left[-\frac{2p^2(n - 2n_0 - 1)^2m}{n + m - 2}\right].
				\end{equation*}\end{linenomath}
				Applying a double union bound over choices of $S$ and $B$, noting that by \expref{Lemma}{lem:binom-bound}, the number of choices $\binom{n}{n_0}\binom{n}{m}$ is at most $4^n/n$, so we have
				\begin{linenomath}\begin{equation*}
						\Pro \left(\exists B_0\in \binom{V}{n_0}, \exists S\in \binom{V}{m}, S\subset B_1 \right)
						\ \le \ \frac{4^n}{n} \exp\left(-\frac{2p^2(n - 2n_0 - 1)^2m}{n + m - 2} \right).
				\end{equation*}\end{linenomath}
				Taking the complement event, we get the desired result.
			\end{proof}
			
			\expref{Lemma}{lem:uni-bad-set-bound} is sufficient to prove parts 2 and 3 of \expref{Theorem}{thm:const-advantage}.
			The following lemmas are direct corollaries of \expref{Lemma}{lem:uni-bad-set-bound}.
			
			\begin{lemma}[Part 2 of \expref{Theorem}{thm:const-advantage}]  \label{prop:uni-sqrt-advantage-ex}
				In the election process on $G\sim G(n, p)$,
				with probability at least
				\begin{equation} \label{eq:uni-sqrt-advantage-ex}
					1 - n^{-1}\exp\left(- .025n \right) = 1 - P_2(n),
				\end{equation}
				$G$ is such that any choice of the Blue camp of at most $n/2 - (.8/p)\sqrt{n} = n_1^B$ members shrinks to size at most $.4n = n_2^B$ the next day.
			\end{lemma}
			
			\begin{proof}[Proof of \expref{Lemma}{prop:uni-sqrt-advantage-ex}]
				Let $n_2 \defeq \left\lfloor n/2 - (.8/p)\sqrt{n} \right\rfloor$ and $m \defeq \left\lceil .4n \right\rceil$.
				By \expref{Lemma}{lem:uni-bad-set-bound}, $G$ is such that every Blue set of at most $n_2$ vertices shrinks to size $m - 1$ with probability at least
				\begin{linenomath}\begin{equation*}  \label{eq:uni-sqrt-advantage-ex-temp1}
						\begin{aligned}
							1 - \frac{4^n}{n} \exp \left(- \frac{2p^2(n - 2n_2 - 1)^2m}{n + m -2} \right)
							= 1 - \frac{1}{n} \exp \left[- \left(\frac{2p^2(n - 2n_2 - 1)^2m}{n + m -2} - 2n \log 2\right)\right].
						\end{aligned}
				\end{equation*}\end{linenomath}
				We have
				\begin{linenomath}\begin{equation*}
						\frac{m}{n + m - 2} \ge \frac{m}{n + m} \ge 2/7, \quad \text{ and } \ n - 2n_2 - 1 \ge \frac{1.6\sqrt{n}}{p} - 1.
				\end{equation*}\end{linenomath}
				Therefore can bound the exponent of the RHS of \expeqref{Equation}{eq:uni-sqrt-advantage-ex-temp1} as follows:
				\begin{linenomath}\begin{equation*}
						\begin{aligned}
							& \frac{2p^2(n - 2n_2 - 1)^2m}{n + m -2} - 2n \log 2
							\ \ge \ \frac{4}{7}p^2\left(\frac{1.6\sqrt{n}}{p} - 1 \right) - 2n\log 2 \\
							& \ = \ \left(\frac{10.24}{7} - 2\log 2\right)n - \frac{12.8}{7}p\sqrt{n} + \frac{4}{7}p^2 \ \ge \ .076n - 1.83\sqrt{n}.
						\end{aligned}
				\end{equation*}\end{linenomath}
				The proof is complete.
			\end{proof}
			
			
			
			\begin{lemma}[Part 3 of \expref{Theorem}{thm:const-advantage}] \label{prop:uni-linear-advantage}
				In the election process on $G\sim G(n, p)$,
				with probability at least
				\begin{linenomath}\begin{equation*}
						1 - n^{-1}\exp\left(-.02n(p^3n - 80)\right) = 1 - P_3(n, p),
				\end{equation*}\end{linenomath}
				$G$ is such that any choice of the Blue camp with at most $.4n$ members shrinks to at most $\frac{1}{3}p(n - 1)$ members the next day.
			\end{lemma}
			
			\begin{proof}[Proof of \expref{Lemma}{prop:uni-linear-advantage}]
				Let $n_2 \defeq \left\lfloor .4n \right\rfloor$ and $m \defeq \lceil \frac{1}{3}p(n-1) \rceil$.
				By \expref{Lemma}{lem:uni-bad-set-bound}, $G$ satisfies that every Blue set of at most $n_2$ vertices shrinks to size $m - 1$ with probability at least
				\begin{equation} \label{eq:uni-linear-advantage-temp1}
					\begin{aligned}
						1 - \frac{4^n}{n} \exp\left(-\frac{2p^2(n - 2n_2 - 1)^2m}{n + m - 2} \right)
						= 1 - \frac{1}{n} \exp\left[-\Bigl(\frac{2p^2(n - 2n_2 - 1)^2m}{n + m - 2} - 2n \log 2\Bigr)\right].
					\end{aligned}
				\end{equation}
				Since $n_2\le .4n$, we have $(n - 2n_2 - 1)^2 \ge (.2n - 1)^2$.
				Furthermore,  $m\ge \frac{1}{3}p(n-1)$ so
				\begin{linenomath}\begin{equation*}
						\frac{m}{n + m - 2}
						\ge \frac{p(n-1)/3}{n + p(n-1)/3 - 1} = \frac{p}{3 + p} \ge \frac{p}{4}.
				\end{equation*}\end{linenomath}
				Therefore we can bound the exponent in the RHS of \expeqref{Equation}{eq:uni-linear-advantage-temp1} as follows:
				\begin{linenomath}\begin{equation*}
						\begin{aligned}
							& \frac{2p^2(n - 2n_2 - 1)^2m}{n + m - 2} - 2n \log 2
							~\ge~  \frac{2p^2(.2n - 1)^2 m}{n + m - 2} - 2n \log 2
							~\ge~ \frac{p^3(.2n - 1)^2}{2} - 2n \log 2 \\
							& ~=~ p^3\left(.02n^2 - .2n + \frac{1}{2}\right) - 2n \log 2
							~\ge~ .02 p^3 n^2 - (.2p^3 + 2\log 2)n + \frac{p^3}{2} \\
							& ~=~ .02 n (p^3n - 10p^3 - 100\log 2)
							~\ge~ .02 n (p^3n - 80),
						\end{aligned}
				\end{equation*}\end{linenomath}
				%
						%
							%
							%
							%
							%
							%
						%
				where the last inequality is due to $10p^3 + 100\log 2 \le 10 + 100\log 2 < 80$.
				The result follows.
			\end{proof}
			
			%
			%
			%
			These lemmas and \expref{Lemma}{prop:uni-large-linear-advantage}, together with \expref{Lemma}{lem:const-advantage-day-1} form the complete ``chain of shrinking'' for the number of Blue vertices to reach $0$ in four days, hence wrapping up the proof of \expref{Theorem}{thm:const-advantage}.
			
			
			
			%
			%
			%
			%
			%
			%
			
			\section{Conclusion} \label{sec:related}
			
			The majority dynamics scheme on a network is a process where each of the $n$ individuals is assigned an initial color, which changes daily to match the majority among their neighbors.
			%
			%
			Our main result, \expref{Theorem}{thm:const-advantage}, when the network is a $G(n, p)$ random graph, yields an explicit %
			lower bound  %
			based on $n, p$ and $c$ for the probability that the side with the initial majority wins. It has two important implications. The first is a surprising phenomenon, which we call
			the  %
			\emph{Power of Few} phenomenon (\expref{Theorem}{thm:const-advantage-ex}), which shows that when $p=1/2$ and $\eps =.1$, $c$ can be set to just $5$, meaning five extra people is all it takes to win a large election with overwhelming odds. The second is an asymptotic dependency between the $\eps, n$ and $c$ (\expref{Theorem}{thm:unbounded-advantage}), which shows that for any fixed $p$, there is a constant $K(p)$ such that choosing $n$ and $c$ both large enough so that $K(p)\max \{n^{-1}, c^{-2} \} < \eps$ will ensure that the winning probability is at least $1 - \eps$.
			
			%
			
			Although the results in this paper only 
			%
			apply to  %
			dense $G(n, p)$ graphs, we do cover sparse graphs in a separate in-progress paper \cite{TranVusparse}, where we obtain the Power of Few phenomenon for $p = \Omega((\log n)/n)$, and discuss the end result (other than a win) for lower values of $p$. We nevertheless %
			mention  %
			one of the main 
			%
			results proved in  %
			the upcoming paper (\expref{Theorem}{thm:general-sparse}), 
			and %
			use it  %
			to prove the main theorem
			%
			of  %
			the paper \cite{fountoulakis} by Fountoulakis, Kang and Makai in \expref{Appendix}{appx:FKM-proof}.
			
			%
			%
			%
			
			
			%
			
			\appendix
			
			%
			%
			%
			%
			%
			%
			%
			
			%
			\section{Proof of Fountoulakis et al's Theorem from \expref{Theorem}{thm:general-sparse}} 
			\label{appx:FKM-proof}
			
			%
			%
			%
			%
			\iffalse %
			\begin{proof}[Proof of \expref{Lemma}{lem:binom-bound}]
				Firstly we show that for a given $n$, $\binom{n}{k}$ is maximized if and only if $k \in \{\lfloor n/2 \rfloor, \lceil n/2 \rceil\}$.
				
				\medskip
				
				\noindent If $k \le \lfloor n/2 \rfloor - 1$, we have $2k \le n - 2$, so $k + 1 < n - k$, thus
				$\displaystyle \binom{n}{k} = \frac{k + 1}{n - k} \binom{n}{k + 1} < \binom{n}{k + 1}$.
				
				\medskip
				
				\noindent If $k \ge \lceil n/2 \rceil + 1$, we have $2k \ge n + 2$, so $n - k + 1 < k$, thus
				$\displaystyle \binom{n}{k} = \frac{n - k + 1}{k} \binom{n}{k - 1} \le \binom{n}{k - 1}$.
				
				\medskip
				
				Therefore if $k$ is not one of $\lfloor n/2 \rfloor$ and $\lceil n/2 \rceil$ (which are either the same or two consecutive numbers), $\binom{n}{k}$ is not maximized.
				Since $\lfloor n/2 \rfloor + \lceil n/2 \rceil = n$, both values for $k$
				%
				maximize $\binom{n}{k}$.  %
				
				It suffices to prove $\binom{n}{\lfloor n/2 \rfloor} < \frac{2^n}{\sqrt{n}}$ for $n \ge 1$.
				%
				For even $m$, this is stated in \cite[2.6.2 Proposition]{invitation}.
				%
				First consider the case $n = 2m$ for $m\in \mathbb{N}$. We have
				\begin{linenomath}\begin{equation*}
						\begin{aligned}
							& \binom{2m}{m} < \frac{2^{2m}}{\sqrt{2m}}
							~\iff~ \frac{\prod_{i=1}^{m-1} (2i + 1) \cdot \prod_{i=1}^m (2i)}{m!\cdot m!} < \frac{2^{2m}}{\sqrt{2m}}
							~\iff~ \frac{\prod_{i=1}^{m-1} (2i + 1)}{m!} < \frac{2^m}{\sqrt{2m}} \\
							& ~\iff~ \frac{1}{m} \prod_{i=1}^{m-1} \frac{2i + 1}{i} < \frac{2^m}{\sqrt{2m}}
							~\iff~ \frac{1}{m}\prod_{i=1}^{m-1} \left(\frac{2i + 1}{i}\right)^2 < \frac{2^{2m}}{2m}
							~\iff~ \prod_{i=1}^{m-1} \left(\frac{2i + 1}{i}\right)^2 < 2^{2m - 1}m.
						\end{aligned}
				\end{equation*}\end{linenomath}
				To proceed, we show that $\left(\frac{2i + 1}{i}\right)^2 < \frac{4i}{i - 1}$ for $i \ge 2$.
				We have
				\begin{linenomath}\begin{equation*}
						\left(\frac{2i + 1}{i}\right)^2 < \frac{4i}{i - 1}
						~\iff~ 4 + \frac{4}{i} + \frac{1}{i^2} \le 4 + \frac{4}{i - 1}
						~\iff~ \frac{1}{i^2} < \frac{4}{i} - \frac{4}{i - 1}
						~\iff~ \frac{1}{i^2} < \frac{4}{i(i - 1)}.
				\end{equation*}\end{linenomath}
				The last inequality is clearly true for $i\ge 2$.
				Suppose first that $m \ge 3$.
				It now suffices to show %
				any of the following equaivalent statements:  %
				\begin{linenomath}\begin{equation*}
						\left(2 + \frac{1}{1}\right)^2 \left(2 + \frac{1}{2}\right)^2 \left[\prod_{i=3}^{m-1} \frac{4i}{i - 1}\right] < 2^{2m - 1}m
						~\iff~ \frac{225}{4}\cdot 4^{m-3}\frac{m-1}{2} < 2^{2m - 1}m
						~\iff~ 225(m - 1) < 256m.
				\end{equation*}\end{linenomath}
				The last inequality clearly holds, so we have shown $\binom{2m}{m} < \frac{2^{2m}}{\sqrt{2m}}$ for $m\ge 3$.
				The cases that $m = 1, 2$ can be verified easily.
				Now consider $n = 2m - 1$ for $m \ge 1$.
				We have
				\begin{linenomath}\begin{equation*}
						\binom{2m - 1}{m - 1} = \frac{1}{2} \binom{2m}{m} < \frac{1}{2} \frac{2^{2m}}{\sqrt{2m}}
						= \frac{2^{2m - 1}}{\sqrt{2m}} < \frac{2^{2m - 1}}{\sqrt{2m - 1}}.
				\end{equation*}\end{linenomath}
				The proof is complete.
			\end{proof}
			\fi  %
			
			%
			
			%
			%
			%
			
			
			%
			%
			%
			%
			In this Appendix we   %
			give a proof of the main theorem by Fountoulakis et al in \cite{fountoulakis} (\expref{Theorem}{FKM}) using the main theorem in our upcoming paper (\expref{Theorem}{thm:general-sparse}).
			
			\begin{proof}
				Assume \expref{Theorem}{thm:general-sparse}. Let $R_0$ and $B_0$ respectively be the initial Red and Blue camps.
				Fix  a constant $0 < c' \le \eps/6$. $\card{R_0}\sim \Bin(n, 1/2)$ since it is a sum of $\Bin(1, 1/2)$ variables. An application of the Berry--Esseen theorem (\expref{Theorem}{thm:esseen}; with $.56 =.56$) implies that
				\begin{linenomath}\begin{equation*}
						\Pr \left(\card{R_0} - \frac{n}{2} \le c'\sqrt{n} \right) \le \Phi(2c') + \frac{.56}{\sqrt{n}},
						\quad \text{ and } \
						\Pr \left(\card{R_0} - \frac{n}{2} \le -c'\sqrt{n} \right) \ge \Phi(-2c') - \frac{.56}{\sqrt{n}},
				\end{equation*}\end{linenomath}
				Thus
				\begin{linenomath}\begin{equation*}
						\begin{aligned}
							& \Pr \left(\abs{\card{R_0} - \frac{n}{2}} \le c'\sqrt{n} \right)
							\ \le \ \left(\Phi(2c') + \frac{.56}{\sqrt{n}}\right)
							- \left(\Phi(-2c') - \frac{.56}{\sqrt{n}}\right) \\
							& \ \le \ \Phi(-2c', 2c') + \frac{20.56}{\sqrt{n}}
							\ \le \  \frac{4c'}{\sqrt{2\pi}} + \frac{20.56}{\sqrt{n}}
							\ \le \ \frac{\eps}{3} + \frac{20.56}{\sqrt{n}} \le \eps/2,
						\end{aligned}
				\end{equation*}\end{linenomath} for sufficiently large $n$.
				
				
				On the other hand, if $\abs{\card{R_0} - n/2} > c'\sqrt{n}$, then one of the sides has more than $n/2 + c' \sqrt{n}$ initial members, which we call the \emph{majority side}.
				Now we apply \expref{Theorem}{thm:general-sparse} with $\eps$ replaced by $\eps/2$.
				Notice that in the setting of \expref{Theorem}{FKM} %
				, %
				if  we have $p =\lambda n^{-1/2} $ for $\lambda$ sufficiently large, then $c' \sqrt n \ge c/ p$, where $c$ is the constant in \expref{Theorem}{thm:general-sparse}. Thus, by  this theorem,  the probability for the majority side to win is at least  $1-\eps/2$, and we are done by the union bound.
			\end{proof}
			
			
			\bibliographystyle{tocplain}   %
			
			%
			%
			%
			%
			\bibliography{bibstrings,v019a006,bibtail}
			
			%
			%
			%
			%
			%
			\begin{tocauthors}
				\begin{tocinfo}[tran]
					Linh Tran\\
					Graduate Student\\
					Department of Mathematics\\
					Yale University\\
					New Haven, CT, USA\\
					l.tran\tocat{}yale\tocdot{}edu \\   %
					\url{https://math.yale.edu/people/linh-tran}      %
				\end{tocinfo}
				\begin{tocinfo}[vu]
					Van Vu\\
					Percey F. Smith Professor of Mathematics and Professor of Statistics \& Data Science\\
					Yale University\\
					New Haven, CT, USA\\
					van.vu\tocat{}yale\tocdot{}edu \\   %
					\url{https://campuspress.yale.edu/vanvu/}      %
				\end{tocinfo}
			\end{tocauthors}
			
			%
%
%
%
%
%
%
%
%
%
%
%
%
%
%
%
%
%
%
%
\begin{tocaboutauthors}
\begin{tocabout}[tran]  %
\textsc{Linh Tran}
is currently a fifth year \phd\ student in the Math department
at Yale University, working under the supervision of Professor Van Vu.
He went to special math classes for gifted children at the high school
division of the Vietnam National University in Ho Chi Minh City.
Encouraged and nurtured by his high school teachers, Tran found his 
passion for Math and was determined to follow this path,
especially after having had little success in other subjects.  
Under the guidance of his advisor, Prof. Van
Vu, he does research in random graph theory, particularly majority dynamics
and generalizations, and random matrix theory, particularly its application
in data analysis and machine learning methods. He is also interested in
combinatorics problems involving the probabilistic method, applications of
geometry in combinatorics, \eg, random fractals, and theoretical computer
science.
\end{tocabout}
\begin{tocabout}[vu]
\textsc{Van Vu} is the P. F. Smith Professor of Mathematics and Professor of 
Data Science at Yale University. Vu was born in Hanoi (Vietnam) in 1970.
He went to special math classes for gifted children at Chu Van An and 
Hanoi-Amsterdam high schools.
In 1987, he went to Hungary for his undergraduate studies, and obtained his
M.\,Sc. in mathematics at %
Eötvös University, Budapest, in 1994.
He received his \phd\ at Yale University, in 1998 under the direction
of László Lovász. 
His research interest includes random structures (random matrices,
random graphs, random functions), probabilistic and additive combinatorics,
and randomized algorithms.
Recently, he %
he has also been  %
interested in real-life applications of data science
(such as NLP and computer vision). 
\end{tocabout}
\end{tocaboutauthors}
			
\end{document}
		
